\section{Theory of Change}

The scaffold is based on a simple theory of change: machine-learning studies often fail scientifically because important research decisions are made out of sequence, remain undocumented, or are disconnected from the final manuscript claim. A model can produce strong performance metrics while the study remains vulnerable to weak question framing, unclear target definition, data leakage, inadequate validation, incomplete ethics review, or unsupported generalization. These are not only technical problems. They are workflow problems.

The proposed scaffold intervenes by imposing a staged sequence from research query to publication. The sequence begins with query framing and study protocol development, then proceeds through data audit, feature planning, modeling plan, validation plan, result bundling, ethics review, reproducibility packaging, manuscript drafting, and submission preparation. Each stage produces an artifact that can be inspected, revised, versioned, and linked to the manuscript.

The scaffold is expected to improve machine-learning scientific practice through five mechanisms. Sequencing keeps the research question and protocol ahead of model selection. Visibility turns hidden assumptions into reviewable artifacts. Traceability links manuscript claims to evidence sources and validation checks. Reviewability gives collaborators and reviewers a clearer path from design decisions to final claims. Reproducibility preserves the materials needed to inspect, rerun, or extend the study.

The expected result is not simply a cleaner repository. The expected result is a manuscript whose claims are better aligned with its evidence chain. Under this model, reproducibility is built during the research process rather than assembled after drafting. Limitations are identified through artifact review rather than added as generic closing language. Validation is designed to support specific claims rather than to report metrics in isolation.

This theory of change also defines the scaffold's boundary conditions. The framework is most useful for empirical machine-learning studies intended for scholarly publication. It is less central for purely exploratory notebooks, production-only engineering systems, or theoretical machine-learning papers where the primary contribution is algorithmic or mathematical. Its main value is in research settings where the credibility of the final claim depends on the integrity of the workflow that produced it.
